\documentclass[12pt]{article}
\usepackage[utf8]{inputenc}
\usepackage{minted}
\usepackage{multirow}
\usepackage{listings}

\usepackage{amssymb,amsmath,amsthm,amsfonts}
\usepackage{calc}
\usepackage{graphicx}
\usepackage{subfigure}
\usepackage{indentfirst}
\usepackage{titlesec}
\newcommand{\sectionbreak}{\clearpage}
\DeclareGraphicsExtensions{.bmp,.png,.pdf,.jpg}
\usepackage{gensymb}

\usepackage{url}
\usepackage{amsmath}
\usepackage{graphicx}
\graphicspath{{images/}}
\usepackage{parskip}
\usepackage{fancyhdr}
\usepackage{vmargin}
\setmarginsrb{2 cm}{2 cm}{2 cm}{2 cm}{1 cm}{1.5 cm}{1 cm}{1.5 cm}

\usepackage[spanish]{babel}
\usepackage[colorlinks=true, allcolors=blue]{hyperref}
\hypersetup{
    colorlinks=true,% make the links colored
}
\usepackage[nolist]{acronym}
\usepackage[table]{xcolor}
\usepackage{url}

\addto\captionsspanish{\renewcommand{\refname}{Bibliografía}}

\begin{document}
\begin{titlepage}

\title{     \textbf{Propuesta de Trabajo Profesional \\ de Ingeniería en Informática}\\[2.5ex]
\textit{Título del Trabajo \\(Segundo renglón)} }

\author{    
             \textbf{Tutor:} Gral. José de San Martín \\[2.5ex]
             \textbf{Co-Tutor:} Gral. Manuel Belgrano \\[2.5ex]
             \textbf{Alumnos} \\[2.5ex]  
             Luis A. Huergo, \textit{(Padrón \# 000.001)}\\ \texttt{ luis.huergo@fi.uba.ar } \\[2.5ex]
             Luis A. Huergo, \textit{(Padrón \# 000.001)}\\ \texttt{ luis.huergo@fi.uba.ar }\\[2.5ex]
             Luis A. Huergo, \textit{(Padrón \# 000.001)}\\ \texttt{ luis.huergo@fi.uba.ar }\\[2.5ex] 
                                                       \\[2.5ex]
            \normalsize{Facultad de Ingeniería, Universidad de Buenos Aires}        \\
       }
\date{}

\end{titlepage}

\maketitle
\thispagestyle{empty} 

\maketitle

{
  \hypersetup{linkcolor=black}
  \tableofcontents
}

%##########################
% Abs
%##########################

\section{Resumen}

Esta sección tiene el objetivo de dar una idea general del trabajo a realizar, cubriendo los siguientes aspectos:
\begin{itemize}
    \item Dificultad técnica
    \item Relevancia de la solución
    \item Aspectos interdisciplinarios
    \item Interacciones de negocio
\end{itemize}

{\bf Esta sección deberá tener aproximadamente  20 renglones.}

%##########################
% ESTADO DE SITUACION
%##########################

\section{Estado de situación}

Descripción y motivación del problema a resolver. Listar los requerimientos de la solución.

\bigskip

%##########################
% OBJETIVOS
%##########################

\section{Objetivos}

El objetivo general del trabajo es ...

Los objetivos particulares son:

\begin{enumerate}
    \item 1
    \item 2
    \item 3
    \begin{enumerate}
        \item 3.1
        \item 3.2
    \end{enumerate}
\end{enumerate}

\bigskip

%##########################
% Esquemas
%##########################

\section{Arquitectura (en caso de ser necesario)}

Incluir diagramas y descripción de la arquitectura en caso que les parezca conveniente.

\bigskip

%##########################
% RECURSOS
%##########################

\section{Tecnologías a utilizar}


\subsection{Software}

\subsection{Hardware (en caso de ser necesario)}

\subsection{Otras herramientas (en caso de ser necesario)}

\bigskip



%##########################
% METODOLOGIA
%##########################

\section{Planificación de trabajo}

Explicar la metodología a utilizar, división del trabajo entre los integrantes, tareas, cronograma, carga horaria.

Cantidad total de horas: xxx hs

\bigskip

\begin{center}
\def\arraystretch{1.5}
\begin{tabular}{ |l|c|c|c|c|c|c|c|c|c|c|c|c| }

 \hline
 \multirow{2}{1em}{Tareas} & \multicolumn{12}{|c|}{Meses}\\  \cline{2-13} &
     1 & 2 & 3 & 4 & 5 & 6 & 7 & 8 & 9 & 10 & 11 & 12 \\  \hline
 Tarea A & \cellcolor{gray} & \cellcolor{gray} &  &  &  &  &  &  &  &  &  &   \\ 
 \hline
Tarea B &  & \cellcolor{gray} & \cellcolor{gray} &  &  &  &  &  &  &  &  &   \\ 
 \hline
Tarea C &  & \cellcolor{gray} & \cellcolor{gray} & \cellcolor{gray} & \cellcolor{gray} & \cellcolor{gray} & &  &  &  &  &   \\ 
 \hline
Tarea D &  &  &  &  &  &  &  &  &  & &  &   \\ 
 \hline
Tarea E &  &  &  &  &  &  &  & & & \cellcolor{gray} & \cellcolor{gray} & \cellcolor{gray}    \\ 
 \hline
\end{tabular}
\end{center}

\bigskip

%\begin{itemize}
%    \item Tarea A [100 horas]: (descripción)
%    \item Tarea B [180 horas]: 
%    \item ...
%\end{itemize}

%\bigskip


%\clearpage


\begin{thebibliography}{99}

    \bibitem{Apache Hive}  Apache Hive TM. URL: \url{hive.apache.org}.

\end{thebibliography}

\end{document}
